% Generated from validated Article Draft Result. Do not hand-edit; revise the structured draft instead.
\section{Open-weightの再加速 — 「Open model」を一つの軸で語らない}
\label{pkg:open}
\noindent\textbf{2025年前半のopen-weight系は、vision-language、portable model、multimodal MoE、thinking/non-thinkingと異なる設計を広げた。同時にhosted APIも増え、「weightsが得られること」と「serviceとして使えること」は別の選択肢になった。}\autocite{src-05f285e9757b,src-1bb16d7254b1,src-3e4a3f3b3d83,src-dc21f366b638}

1月のQwen2.5-VLは3B・7B・72Bのvision-language familyとして公開され、4月末のQwen3ではdenseとMoE、thinking/non-thinking modesを含む新系列へ進んだ。同じQwenでも、H1の技術的焦点は視覚理解からreasoning/deploymentの組合せへ広がった。\autocite{src-3e4a3f3b3d83,src-dc21f366b638}


Qwen2.5-Maxは1月末に大規模MoEとして発表され、Alibaba Cloud API availabilityも示された。ここで重要なのは、hosted serviceとしての利用可能性と、weight distributionの性質を同じ「open」という語でまとめないことである。\autocite{src-8c485019ad8d}


3月のGemma 3は1B・4B・12B・27Bを含むportable open-model familyとして登場した。model sizeの幅はdeployment choiceを広げるが、weights availabilityだけからtraining data transparencyまで推論してはならない。\autocite{src-1bb16d7254b1}


4月のLlama 4 Scout / Maverickはnatively multimodalなMoE modelとして公開され、同月末にはLlama API limited previewが別に提示された。model weight releaseとhosted API previewは異なるproduct identityである。\autocite{src-05f285e9757b,src-fca961b24baf}


H1のopen-weight競争をparameter countだけで比較すると、MoE、multimodality、thinking mode、device portability、hosted APIという実用上の差を落とす。採用時にはlicense、weights、runtime、service availability、data disclosureを分解して評価する必要がある。\autocite{src-05f285e9757b,src-1bb16d7254b1,src-3e4a3f3b3d83}


\begin{claimboundary}
本章で「open-weight」と呼ぶ場合、weightsの利用可能性を中心に記述し、license条件やtraining/data transparencyが同等であることを意味しない。性能比較も各社資料の主張を独立測定値へ読み替えない。
\end{claimboundary}
